\documentclass[aspectratio=169]{beamer}

\mode<presentation>
{
 \usetheme[reversetitle,notitle,noauthor]{Wien}
}

\usepackage{url}
\usepackage{graphicx}
\graphicspath{{./}{./Figures/}}

\usepackage{appendixnumberbeamer}
\usepackage{amsmath}
\usepackage{booktabs}
\usepackage{array}
\usepackage{xcolor}

% To avoid a warning from the hyperref package:
\pdfstringdefDisableCommands{%
  \def\translate{}%
}

\definecolor{tuwred}{RGB}{140,0,0}
\definecolor{darkgreen}{RGB}{0,120,60}
\definecolor{darkblue}{RGB}{0,70,140}
\definecolor{darkorange}{RGB}{190,100,0}

\newcommand{\good}[1]{\textcolor{darkgreen}{\textbf{#1}}}
\newcommand{\warn}[1]{\textcolor{darkorange}{\textbf{#1}}}
\newcommand{\bad}[1]{\textcolor{tuwred}{\textbf{#1}}}

%%%%%%%%%%%%%%%%%%%%%%%%%%%%%%%%%%%%%%%%%%%%%%%%%%%%%%%%%%%%%%%%%%%%%%%%%%%%%
%%%%%%%%%%%%%%%%%%%%%%%%%%%%%%%%%%%%%%%%%%%%%%%%%%%%%%%%%%%%%%%%%%%%%%%%%%%%%
\title[EvoPress Extension]{Extending EvoPress with Joint Multi-Method Compression}
\subtitle{Supervisor Meeting: Experimental Progress and Thesis Direction}
\author[Kerim Halilovi\'c]{Kerim Halilovi\'c}
\institute[TU Wien]{TU Wien, Vienna, Austria}
\date{\today}

\begin{document}

\begin{frame}
  \titlepage
\end{frame}

%%%%%%%%%%%%%%%%%%%%%%%%%%%%%%%%%%%%%%%%%%%%%%%%%%%%%%%%%%%%%%%%%%%%%%%%%%%%%

\begin{frame}{Goal of the current thesis direction}
\begin{block}{Research question}
Can an EvoPress-style evolutionary search jointly optimize \textbf{structural pruning} and \textbf{weight quantization} better than applying them independently?
\end{block}

\vspace{0.2cm}
\begin{itemize}
    \item Original EvoPress studies depth pruning, sparsity, and quantization mostly as separate applications.
    \item My extension investigates \textbf{combined compression}: depth pruning + sparsity/quantization.
    \item Main focus so far: whether joint search discovers structures that remain robust after quantization.
\end{itemize}

\vspace{0.25cm}
\begin{block}{Current status}
Pipeline is implemented and reproducible; first broad results show \textbf{mask-compatibility effects}, but not yet a general proof of joint-search superiority.
\end{block}
\end{frame}

%%%%%%%%%%%%%%%%%%%%%%%%%%%%%%%%%%%%%%%%%%%%%%%%%%%%%%%%%%%%%%%%%%%%%%%%%%%%%

\begin{frame}{Experimental setup and practical constraint}
\begin{columns}[T]
\begin{column}{0.52\textwidth}
\textbf{Models}
\begin{itemize}
    \item \texttt{Mistral-7B-v0.3}: depth-pruning validation
    \item \texttt{TinyLlama-1.1B-Chat}: sparse/quant and joint experiments
\end{itemize}

\vspace{0.2cm}
\textbf{Main metric}
\begin{itemize}
    \item WikiText2 perplexity (PPL), lower is better
    \item KL divergence used as search fitness
\end{itemize}
\end{column}

\begin{column}{0.44\textwidth}
\textbf{Hardware constraint}
\begin{itemize}
    \item GPU varies: Tesla T4 or A40
    \item A40 provides $\sim$45 GB VRAM
    \item Container RAM limit: \bad{16 GB CPU RAM}
\end{itemize}

\vspace{0.2cm}
\begin{block}{Consequence}
Full Mistral sparse database generation is blocked by CPU RAM, so combined-method testing moved to TinyLlama.
\end{block}
\end{column}
\end{columns}
\end{frame}

%%%%%%%%%%%%%%%%%%%%%%%%%%%%%%%%%%%%%%%%%%%%%%%%%%%%%%%%%%%%%%%%%%%%%%%%%%%%%

%%%%%%%%%%%%%%%%%%%%%%%%%%%%%%%%%%%%%%%%%%%%%%%%%%%%%%%%%%%%%%%%%%%%%%%%%%%%%

\begin{frame}{Possible hardware requirements for Mistral-7B experiments}
\begin{columns}[T]
\begin{column}{0.48\textwidth}
\textbf{Minimum pilot}
\begin{itemize}
    \item A40 / L40S GPU, 48 GB VRAM
    \item 64 GB CPU RAM
    \item 150 GB local NVMe storage
    \item 4k--8k calibration tokens
    \item Limited quantization scope first
\end{itemize}
\end{column}

\begin{column}{0.48\textwidth}
\textbf{Recommended full experiment}
\begin{itemize}
    \item A100 / H100 GPU, 80 GB VRAM
    \item 128 GB CPU RAM
    \item 500 GB local NVMe storage
    \item 8k--32k calibration tokens
    \item Broader all-projection quantization
\end{itemize}
\end{column}
\end{columns}

\vspace{0.25cm}
\begin{block}{Practical takeaway}
The GPU is less critical than the current CPU RAM limit. 
A single large GPU with enough host RAM is preferable to multiple smaller GPUs.
\end{block}
\end{frame}

%%%%%%%%%%%%%%%%%%%%%%%%%%%%%%%%%%%%%%%%%%%%%%%%%%%%%%%%%%%%%%%%%%%%%%%%%%%%%

\begin{frame}{Mistral-7B depth pruning: quality curve}
\begin{table}
\centering
\begin{tabular}{lrrrr}
\toprule
Requested depth sparsity & Dropped attn & Dropped MLP & WikiText2 PPL & Train PPL \\
\midrule
Dense & 0 & 0 & \textbf{5.35} & -- \\
12.5\% & 4 & 4 & \good{6.75} & 6.27 \\
25.0\% & 8 & 8 & 14.70 & 13.30 \\
37.5\% & 12 & 12 & 51.69 & 48.20 \\
50.0\% & 16 & 16 & 371.75 & 337.00 \\
\bottomrule
\end{tabular}
\end{table}

\vspace{0.2cm}
\begin{itemize}
    \item Moderate depth pruning is feasible, but quality drops nonlinearly.
    \item 37.5\% and 50\% are difficult operating points.
    \item This establishes the baseline structural-pruning behavior on a 7B model.
\end{itemize}
\end{frame}

%%%%%%%%%%%%%%%%%%%%%%%%%%%%%%%%%%%%%%%%%%%%%%%%%%%%%%%%%%%%%%%%%%%%%%%%%%%%%

\begin{frame}{Why evolutionary search matters: baselines}
\begin{table}
\centering
\begin{tabular}{lrrr}
\toprule
Sparsity & EvoPress & Random median & Late-layer heuristic \\
\midrule
12.5\% & \good{6.75} & 21.50 & 61.16 \\
25.0\% & \good{14.70} & 2490.00 & 541.00 \\
37.5\% & \good{51.69} & 4476.00 & 1419.00 \\
50.0\% & \good{371.75} & 15516.00 & 4992.00 \\
\bottomrule
\end{tabular}
\end{table}

\vspace{0.2cm}
\begin{block}{Finding}
The location of removed modules matters strongly. EvoPress finds nontrivial masks that simple random or late-layer removal misses.
\end{block}
\end{frame}

%%%%%%%%%%%%%%%%%%%%%%%%%%%%%%%%%%%%%%%%%%%%%%%%%%%%%%%%%%%%%%%%%%%%%%%%%%%%%

\begin{frame}{TinyLlama: from feasibility to broader quantization}
\begin{table}
\centering
\begin{tabular}{lrrr}
\toprule
Quantization scope & Modules & Database size & WikiText2 PPL \\
\midrule
Dense TinyLlama & -- & -- & \textbf{8.97} \\
\texttt{q\_proj} only, searched 3-bit & 22 & 529 MB & \good{9.00} \\
All attention projections & 88 & 1.19 GB & database feasible \\
All transformer projections, uniform 3-bit & 154 & 5.55 GB & \good{13.62} \\
All transformer projections, searched 3-bit & 154 & 5.55 GB & 14.58 \\
\bottomrule
\end{tabular}
\end{table}

\vspace{0.2cm}
\begin{block}{Important scope note}
``All transformer projections'' means \texttt{q,k,v,o,gate,up,down}; embeddings, norms, and \texttt{lm\_head} are not included.
\end{block}
\end{frame}

%%%%%%%%%%%%%%%%%%%%%%%%%%%%%%%%%%%%%%%%%%%%%%%%%%%%%%%%%%%%%%%%%%%%%%%%%%%%%

\begin{frame}{Independent vs. joint compression design}
\begin{columns}[T]
\begin{column}{0.48\textwidth}
\begin{block}{Independent composition}
\begin{enumerate}
    \item Run depth search
    \item Run quantization search
    \item Combine both selected configurations
    \item Evaluate final model
\end{enumerate}
\end{block}

\vspace{0.15cm}
\textbf{Risk:} each method is optimized as if the other method were absent.
\end{column}

\begin{column}{0.48\textwidth}
\begin{block}{Joint search}
\begin{enumerate}
    \item One candidate contains both:
    \begin{itemize}
        \item depth mask
        \item bit allocation
    \end{itemize}
    \item Mutate and select under one shared fitness
    \item Enforce active 3-bit budget
\end{enumerate}
\end{block}

\vspace{0.15cm}
\textbf{Goal:} find depth masks that tolerate quantization better.
\end{column}
\end{columns}
\end{frame}

%%%%%%%%%%%%%%%%%%%%%%%%%%%%%%%%%%%%%%%%%%%%%%%%%%%%%%%%%%%%%%%%%%%%%%%%%%%%%

\begin{frame}{Narrow joint experiment: \texttt{q\_proj} quantization}
\begin{table}
\centering
\begin{tabular}{rrrr}
\toprule
Seed & Depth-only PPL & Joint depth+quant PPL & Joint $-$ depth \\
\midrule
0 & 12.02 & \good{11.03} & -0.99 \\
1 & \good{12.28} & 12.55 & +0.27 \\
2 & \good{11.03} & 11.72 & +0.69 \\
\midrule
Mean & 11.78 & 11.77 & -0.01 \\
\bottomrule
\end{tabular}
\end{table}

\vspace{0.2cm}
\begin{itemize}
    \item Joint search worked and replayed correctly.
    \item Added \texttt{q\_proj} quantization had a small average penalty: $\sim$0.15 PPL.
    \item But: no consistent improvement over depth-only across three seeds.
\end{itemize}
\end{frame}

%%%%%%%%%%%%%%%%%%%%%%%%%%%%%%%%%%%%%%%%%%%%%%%%%%%%%%%%%%%%%%%%%%%%%%%%%%%%%

\begin{frame}{Broad joint experiment: all transformer projections}
\begin{table}
\centering
\begin{tabular}{lr}
\toprule
Configuration & WikiText2 PPL \\
\midrule
Dense & \textbf{8.97} \\
Depth-only seed 0 & 12.02 \\
Uniform all-projection 3-bit & 13.62 \\
Searched all-projection 3-bit & 14.58 \\
Depth + uniform all-projection 3-bit & 21.77 \\
Depth + independently searched all-projection 3-bit & 23.81 \\
Joint active-budget depth + all-projection 3-bit & 21.47 \\
Joint depth mask + uniform 3-bit & \good{21.25} \\
\bottomrule
\end{tabular}
\end{table}

\begin{block}{Key result}
The best broad combined model is not the learned mixed-bit allocation; it is the \textbf{jointly discovered depth mask} paired with \textbf{uniform 3-bit quantization}.
\end{block}
\end{frame}

%%%%%%%%%%%%%%%%%%%%%%%%%%%%%%%%%%%%%%%%%%%%%%%%%%%%%%%%%%%%%%%%%%%%%%%%%%%%%

\begin{frame}{Attribution: what actually improved?}
\small
\begin{table}
\centering
\begin{tabular}{p{0.42\textwidth} r p{0.35\textwidth}}
\toprule
\textbf{Comparison} & \textbf{PPL change} & \textbf{Interpretation} \\
\midrule
Joint vs. independent searched composition & -2.34 & Joint search helps \\
Joint vs. uniform composition & -0.30 & Small improvement \\
Joint mask + uniform vs. original mask + uniform & -0.52 & Better quantization-compatible mask \\
Joint allocation vs. uniform on joint mask & +0.22 & Mixed-bit allocation is worse \\
Searched quant-only vs. uniform quant-only & +0.96 & Uniform is stronger baseline \\
\bottomrule
\end{tabular}
\end{table}

\begin{block}{Supported interpretation}
Joint search found a \textbf{depth mask that is more robust to broad quantization}. It did \textbf{not} find a better nonuniform bit allocation.
\end{block}
\end{frame}

%%%%%%%%%%%%%%%%%%%%%%%%%%%%%%%%%%%%%%%%%%%%%%%%%%%%%%%%%%%%%%%%%%%%%%%%%%%%%

\begin{frame}{Proposed thesis contribution}
\begin{block}{Working thesis framing}
Extending EvoPress to \textbf{joint multi-method compression}: analyzing when structural pruning decisions should be optimized together with weight-level quantization.
\end{block}

\vspace{0.2cm}
\textbf{Current evidence supports:}
\begin{itemize}
    \item Joint search is technically feasible and replayable.
    \item Broad transformer-projection quantization can be combined with depth pruning.
    \item Joint search can discover a depth mask that tolerates quantization better.
    \item Uniform bit allocation is currently stronger than learned mixed-bit allocation.
    \item The contribution is not ``joint search always wins.''
\end{itemize}
\end{frame}

%%%%%%%%%%%%%%%%%%%%%%%%%%%%%%%%%%%%%%%%%%%%%%%%%%%%%%%%%%%%%%%%%%%%%%%%%%%%%

\begin{frame}{Limitations and next steps}
\begin{columns}[T]
\begin{column}{0.48\textwidth}
\textbf{Current limitations}
\begin{itemize}
    \item Broad all-projection joint result is currently one seed.
    \item Evaluation is mostly WikiText2 PPL.
    \item No final effective-size or latency measurements yet.
\end{itemize}
\end{column}

\begin{column}{0.48\textwidth}
\textbf{Proposed next step}
\begin{itemize}
    \item Run a stronger independent vs. joint comparison.
    \item Use larger search budgets: more generations, offspring, and calibration tokens.
    \item Repeat the comparison across multiple seeds.
\end{itemize}
\end{column}
\end{columns}

\vspace{0.25cm}
\begin{block}{Goal}
Validate whether joint search consistently discovers depth masks that are more robust to quantization.
\end{block}
\end{frame}

%%%%%%%%%%%%%%%%%%%%%%%%%%%%%%%%%%%%%%%%%%%%%%%%%%%%%%%%%%%%%%%%%%%%%%%%%%%%%%%%%%%%%%%%%%%%%%%%%

\end{document}
